% Library CV analytical_fr — transcription of writing-samples/cvs/Moumen_Nermine_5.pdf (CV12, 2026-09-21)
\documentclass[10pt]{article}
\usepackage[letter]{nermine-cv}
\begin{document}
\cvheader{Moumen Nermine}{Étudiante en dernière année du cycle ingénieur en Chimie Industrielle}{nermine.moumen@insat.ucar.tn}{+21627040414}{linkedin.com/in/nermine-moumen}

\section{FORMATION}
\entry{Institut National des Sciences Appliquées et de Technologie (INSAT)}{2022–Présent}
\summary{Étudiante en Chimie Industrielle, spécialité Procédés et Environnement, avec une formation solide en chimie analytique et techniques de caractérisation, complétée par des bases en science des matériaux, électrochimie, génie des procédés et HSE en laboratoire. Orientée vers le développement, la validation et l'optimisation de méthodes d'analyse chimique. À la recherche d'un PFE.}

\section{EXPÉRIENCE PROFESSIONNELLE}
\entry{SEGULA Technologies - Stage en Modélisation Électrochimique des batteries}{06/2026–08/2026}
\begin{bul}
\item Développement d'un modèle numérique (Single Particle Model, \textbf{MATLAB}) d'une cellule Li-ion, résolu par FVM.
\item Optimisation des paramètres du modèle, réduisant l'erreur de prédiction de \textbf{plus de 80 \%} par rapport aux données expérimentales.
\item Validation des résultats par comparaison à des essais de décharge et au logiciel \textbf{COMSOL Multiphysics}.
\end{bul}
\entry{Control Union - Stage en Contrôle Qualité}{08/2025–09/2025}
\begin{bul}
\item Analyse de lubrifiants (viscosité et métaux lourds par \textbf{ICP-OES}).
\item Préparation et interprétation de rapports d'analyse pour vérifier la conformité aux normes \textbf{ASTM D5185} et \textbf{ISO 17025}.
\end{bul}
\entry{AGIL - Stage en Contrôle Qualité}{07/2025–08/2025}
\begin{bul}
\item Réalisation d'analyses physico-chimiques (\textbf{point d'éclair} et \textbf{distillation atmosphérique}) de carburants diesel et Jet A-1.
\end{bul}
\entry{ETAP - Stage en Caractérisation d'hydrocarbures et traitement des Eaux industrielles}{07/2024–08/2024}
\begin{bul}
\item Formation aux techniques de géochimie organique (GC, GC-MS, pyrolyse Rock-Eval) pour la caractérisation d'hydrocarbures.
\item Analyses physico-chimiques des eaux (pH, salinité, bicarbonates, sulfates, chlorures) par \textbf{chromatographie ionique} et \textbf{ICP-AES}.
\item Étude des méthodes de traitement des eaux pétrolières par échange d'ions et par coagulation-floculation.
\end{bul}
\section{EXPÉRIENCE COMPLÉMENTAIRE}
\entry{Cascade - Enseignante de Chimie en Ligne (Temps partiel)}{2024–Présent}
\begin{bul}
\item Accompagnement individualisé en chimie pour des élèves de Terminale française, couvrant les concepts du niveau examen.
\end{bul}

\section{PROJETS ACADÉMIQUES}
\project{Validation de Méthodes de Dosage du Sodium et du Potassium par Spectrométrie d'Émission de Flamme}
\begin{bul}
\item Validation de \textbf{trois protocoles de quantification} (étalonnage externe, étalonnage interne, ajouts dosés) pour le dosage de Na$^+$ et K$^+$, selon la norme \textbf{NT ISO 9964-3}.
\item Correction des interférences d'ionisation par étalon interne et évaluation de l'effet de matrice via le taux de recouvrement.
\end{bul}
\project{Polymères à Impression Moléculaire (MIPs) en Applications Pharmaceutiques}
\begin{bul}
\item Synthèse et caractérisation d'un MIP sélectif au sulpiride par polymérisation radicalaire, avec extraction Soxhlet, suivi par spectrophotométrie UV-Vis et validation par cinétique de libération contrôlée (240 min).
\item Synthèse d'un MIP core-shell MOF/polymère pour la délivrance ciblée de doxorubicine, avec \textbf{85 \% de taux d'encapsulation}.
\end{bul}
\project{Prédiction des Équilibres Liquide–Liquide pour Systèmes Ternaires d'Extraction}
\begin{bul}
\item Formulation et résolution du système thermodynamique LLE (Rachford-Rice, continuation par longueur d'arc).
\item Réimplémentation multi-langages (Python, MATLAB, Mathematica) avec équivalence démontrée entre méthodes Hunter-Nash et McCabe-Thiele.
\end{bul}

\section{COMPÉTENCES TECHNIQUES}
\begin{bul}
\item \textbf{Techniques de Laboratoire :} Chromatographie (HPLC, GC, CCM), Spectrométrie (MS, ICP-OES), Spectroscopie (RMN, FTIR, UV-Vis), Diffraction (XRD), Microscopie (MEB, MET), Électrochimie (Voltamétrie cyclique, polarographie).
\item \textbf{Simulation :} Aspen Plus et HYSYS, COMSOL Multiphysics\sep\textbf{Programmation :} \textit{MATLAB, Mathematica, Python}
\item \textbf{Données \& Bureautique :} Pack Microsoft Office, OpenLCA\sep\textbf{Normes :} ISO 9001, ISO 14001, ISO 45001, ISO 17025, ISO 14040
\end{bul}

\section{ACTIVITÉS EXTRASCOLAIRES}
\textbf{IEEE INSAT} — Chef d'organisation de TRSYP 3.0 ; Responsable des activités sociales – MENASYP ; Présidente de l'EMBS.\\
\textbf{THEATRO INSAT} — Membre de l'équipe de danse.\par

\section{LANGUES ET CERTIFICATS}
\twocol{Langues}{Anglais (Avancé), Français (Bilingue), Arabe (Natif).}{Certificats}{Aspen HYSYS – MATLAB}
\end{document}
